\begin{example}
\label{example-spec-Zx}
In this example we describe $X = \Spec(\mathbf{Z}[x])$.
Fix $\mathfrak p \in X$.
Let $\phi : \mathbf{Z} \to \mathbf{Z}[x]$ and notice
that $\phi^{-1}(\mathfrak p) \in \Spec(\mathbf{Z})$.
If $\phi^{-1}(\mathfrak p) = (q)$ for $q$ a prime number $q > 0$,
then $\mathfrak p$ corresponds to a prime in $(\mathbf{Z}/(q))[x]$.
The zero prime of that polynomial ring gives $\mathfrak p=(q)$.
Every nonzero prime is generated by an irreducible polynomial
$\overline{f_q}\in(\mathbf{Z}/(q))[x]$.
Choose any coefficient lift $f_q\in\mathbf{Z}[x]$ of the same degree.
The inverse image of $(\overline{f_q})$ is $(q,f_q)$: if the reduction
of $F\in\mathbf Z[x]$ is $\overline A\,\overline{f_q}$, choose a lift
$A\in\mathbf Z[x]$; then $F-Af_q\in q\mathbf Z[x]$.
The converse containment follows by reduction.
Thus the primes over $(q)$ are $(q)$ and the ideals $(q,f_q)$ with
$\overline{f_q}$ irreducible. The lift $f_q$ itself need not be
irreducible in $\mathbf Z[x]$: for $q=5$, the degree-one reduction
of $2x+2$ is irreducible, while $2x+2=2(x+1)$ in $\mathbf Z[x]$.
Now assume that $\phi^{-1}(\mathfrak p)=(0)$.
In this case, if $\mathfrak p = (0)$ there is nothing more to prove. Otherwise, $\mathfrak p$ contains nonconstant polynomials
and it contains no nonzero integer. Factoring any nonzero member
in the unique factorization domain $\mathbf Z[x]$ and using primality
gives an irreducible $f\in\mathfrak p$. It is nonconstant because
$\mathfrak p\cap\mathbf Z=(0)$, and it is primitive because a
nonunit integer content would be a nontrivial factor.
Every other nonconstant irreducible $g\in\mathfrak p$ is likewise primitive.
Gauss' lemma makes $f$ and $g$ irreducible in $\mathbf Q[x]$.
Moreover, if $g=(d/e)f$ with coprime nonzero integers $d,e$,
the equality $eg=df$ has contents $|e|$ and $|d|$, respectively.
Thus $|e|=|d|=1$, so association over $\mathbf Q[x]$ implies
association over $\mathbf Z[x]$ for these original polynomials.
If $f$ and $g$ were not associates, the Euclidean algorithm in
$\mathbf Q[x]$ would give $1=af+bg$ with $a,b\in\mathbf Q[x]$.
Choose a positive common denominator $m$ of every coefficient of
$a$ and $b$, and put $\bar a=ma$, $\bar b=mb$ in $\mathbf Z[x]$.
Then $m=\bar a f+\bar b g\in\mathfrak p\cap\mathbf Z$, a contradiction.
Therefore all irreducibles belonging to $\mathfrak p$ are associates of $f$.
For every nonzero $F\in\mathfrak p$, a factorization of $F$ has
an irreducible factor in $\mathfrak p$ by primality, and that factor
is a unit multiple of $f$. Thus $f$ divides every $F\in\mathfrak p$,
while $f\in\mathfrak p$, proving $\mathfrak p=(f)$.
Conversely, every nonconstant irreducible $f\in\mathbf Z[x]$
generates a prime ideal, because $\mathbf Z[x]$ is a unique
factorization domain. Its contraction to $\mathbf Z$ is zero:
a nonzero constant cannot be divisible by a positive-degree polynomial.
Together with the separately retained zero prime, this describes
all primes with zero contraction.
\end{example}